\chapter{Abortion}
\index{Abortion}

\section*{Definition and Overview}
Abortion, also known as termination of pregnancy, is the medical or surgical ending of a pregnancy before birth. Laws, gestational limits, access arrangements, and funding vary by jurisdiction. It can be provided using medication alone (medical abortion) or by a procedure (surgical abortion). When performed by trained providers in a regulated clinical setting, abortion is generally very safe, with serious complications being uncommon. This chapter explains the methods available, what to expect, the care provided around the procedure, and the situations in which medical help should be sought.

\section*{Epidemiology}
Abortion is common worldwide, and most abortions occur early in pregnancy. Medical abortion accounts for a large share of procedures in many settings. People seek abortion for a wide range of reasons, including contraceptive failure, risks to maternal health, fetal abnormality, and personal, social, or economic circumstances. Cost, distance, stigma, legal restrictions, and availability of trained providers can create barriers to safe care.

\section*{Aetiology and Pathophysiology}
Abortion is not a disease, and its "causes" are the many circumstances in which a woman decides not to continue a pregnancy.

\begin{itemize}
    \item \textbf{Unintended pregnancy:} contraceptive failure, including missed doses or failed methods
    \item \textbf{Maternal health reasons:} conditions in which continuing the pregnancy carries significant risk
    \item \textbf{Fetal abnormality:} a diagnosis made during early pregnancy screening or testing
    \item \textbf{Personal and social circumstances:} including financial, relationship, housing, and family situations
    \item \textbf{Timing:} a pregnancy occurring at a difficult or unplanned point in a woman's life
\end{itemize}

The decision to end a pregnancy is deeply personal, and non-judgemental counselling and support should be available regardless of the reason.

\section*{Clinical Features}
The experience of an abortion varies with the method and gestational age, but some features are common.

\begin{itemize}
    \item Medical abortion involves taking medication that ends the pregnancy, followed by cramping and bleeding over several hours to days
    \item Surgical abortion is a short procedure that removes the pregnancy, with mild cramping and bleeding afterwards
    \item Bleeding is typically similar to a heavy period and settles within one to two weeks
    \item Pain is usually controlled with simple analgesia
    \item Nausea, fatigue, and breast tenderness may persist briefly as pregnancy hormones fall
    \item Most women feel relief; some experience sadness or mixed emotions, which are normal and usually settle with support
\end{itemize}

\section*{Differential Diagnosis}
Before an abortion is provided, the situation must be confirmed and certain conditions excluded.

\begin{itemize}
    \item \textbf{Miscarriage:} a pregnancy that has already ended spontaneously
    \item \textbf{Ectopic pregnancy:} a pregnancy implanted outside the womb, which requires different, urgent management
    \item \textbf{Molar pregnancy:} an abnormal pregnancy with specific follow-up requirements
    \item \textbf{Ongoing pregnancy after medical abortion:} more likely when medication is not taken as directed
    \item \textbf{Incomplete abortion:} retained pregnancy tissue requiring further treatment
\end{itemize}

\section*{When to Seek Medical Care}
After an abortion, the clinic or a doctor should be contacted promptly for heavy bleeding, severe pain, fever, or signs of infection. A follow-up review, usually including ultrasound, is arranged within a few weeks to confirm the abortion is complete.

\textbf{Contact your local emergency services for:}
\begin{itemize}
    \item Very heavy bleeding that soaks pads rapidly or passes large clots
    \item Severe or worsening abdominal pain
    \item Fever, chills, or foul-smelling vaginal discharge
    \item Fainting, light-headedness, or collapse
    \item Chest pain or difficulty breathing
\end{itemize}

\section*{Diagnosis and Investigations}
Assessment before abortion is straightforward and confirms the diagnosis and safety of the chosen method.

\begin{itemize}
    \item \textbf{Confirmation of pregnancy:} a urine or serum pregnancy test
    \item \textbf{Ultrasound:} to date the pregnancy and confirm that it is inside the womb rather than ectopic
    \item \textbf{Blood tests:} haemoglobin, and blood group testing so that Rh-negative women can receive anti-D immunoglobulin to protect future pregnancies
    \item \textbf{Gestational age assessment:} the method offered depends on how many weeks pregnant the woman is
    \item \textbf{Discussion of options, consent, and follow-up,} including contraception for afterwards
\end{itemize}

\section*{Management}
The method of abortion is selected according to gestational age, availability, and the woman's preference.

\begin{itemize}
    \item \textbf{Medical abortion:} mifepristone followed by misoprostol, used in early pregnancy, usually within the first nine weeks
    \item \textbf{Surgical abortion:} suction aspiration, typically performed in the first trimester under local or general anaesthesia
    \item \textbf{Later-term procedures:} for pregnancies beyond the first trimester, surgical management is usually needed, sometimes over one or two days
    \item \textbf{Pain relief and emotional support} during and after the procedure
    \item \textbf{Contraception:} advice and provision of an effective method so that a further abortion is not needed
    \item \textbf{Counselling:} available before and after for those who want it
    \item \textbf{Follow-up:} review to confirm completion of the abortion and normal recovery
\end{itemize}

\section*{Complications}
Serious complications are uncommon when abortion is provided safely.

\begin{itemize}
    \item Bleeding, including incomplete abortion with retained tissue
    \item Infection, which may require antibiotics
    \item Injury to the cervix or, rarely, perforation of the uterus during surgical abortion
    \item Ongoing pregnancy when a medical abortion does not complete
    \item Emotional responses, including sadness, guilt, or regret in some women
    \item Complications are far more frequent with unsafe or unregulated procedures, which is why access to safe, legal care is so important
\end{itemize}

\section*{Prognosis and Outcomes}
Safe abortion is one of the lowest-risk medical procedures, with serious complications being rare. Fertility is preserved, and the great majority of women who wish to conceive subsequently go on to have normal pregnancies. Physical recovery is usually rapid, with bleeding settling within a couple of weeks and normal activities resuming quickly. Most women do not regret their decision, and those who experience distress typically improve with support and time.

\section*{Prevention and Self-Care}
\begin{itemize}
    \item Use effective contraception consistently if pregnancy is not desired
    \item Know about emergency contraception for use after unprotected sex
    \item Plan and attend follow-up after an abortion so that any problem is detected early
    \item Rest, use simple analgesia, and avoid tampons, swimming, or penetrative sex until bleeding has settled
    \item Arrange contraception immediately after an abortion, as fertility returns quickly
    \item Seek counselling or support if feelings are mixed or difficult
    \item Know the emergency warning signs and where to get help, day or night
\end{itemize}

\section*{Sources and Further Reading}
The following organisations provide reliable, non-judgemental information about abortion for women, families, and clinicians.

\begin{itemize}
    \item NHS. \textit{Abortion (termination of pregnancy): overview and what to expect.} https://www.nhs.uk/conditions/abortion/
    \item World Health Organization. \textit{Abortion care and safe abortion guidance.} https://www.who.int/
    \item MSD Manual. \textit{Abortion: professional overview of methods and complications.} https://www.msdmanuals.com/
    \item MedlinePlus. \textit{Abortion health information.} https://medlineplus.gov/abortion.html
\end{itemize}
